% GENERATED FILE. Edit canonical lessons, then run npm run generate:books.
% canonical-source-hash: fnv1a64:d77ef2d1a4410d58
% canonical-lessons: BN-C318-biplob, BN-C318-andolon, BN-C318-songgram, BN-C318-protishtha-kora, BN-C318-mukti

\chapter{Turning Points}
\label{ch:bn-turning-points}
\hlchaptermodality{318}

\begin{chapteropening}
\textbf{By the end of this chapter:} \emph{I can talk about revolutions, movements and what people founded.}

Complete the last lesson of chapter 318: I can talk about revolutions, movements and what people founded.
\end{chapteropening}

\section[\texorpdfstring{\bn{বিপ্লব} (biplôb)}{বিপ্লব (biplôb)}]{\bn{বিপ্লব} (biplôb) — a revolution}

\label{lesson:BN-C318-biplob}



\begin{quote}
Before the new one: say the Bengali for to surrender, then the Bengali for defeat.
\end{quote}

\subsection*{You'll want to know: \bn{বিপ্লব}}
\textbf{\bn{বিপ্লব}} — \emph{biplôb} — \textquotedblleft{}a revolution\textquotedblright{}.

\begin{quote}
\bn{বিপ্লবের} \bn{পরে} \bn{দেশে} \bn{নতুন} \bn{সরকার} \bn{এল।}
\end{quote}

\begin{quote}
\emph{biplôber pôre deshe nôtun sôrkār elô.}
\end{quote}

\begin{quote}
After the revolution a new government came to power.
\end{quote}

\begin{cousinweb}
\textbf{\bn{বিপ্লব}} (\emph{biplôb}) is Sanskrit \emph{viplava}, 'tumult, a revolt; ruin' (Monier-Williams), from \emph{vi-√plu}, 'to float apart'.
\end{cousinweb}

\begin{grammarlens}[title={What you've built}]
One more word for telling what happened.
\end{grammarlens}

\subsection*{Your turn}
\begin{itemize}
\raggedright
  \item \emph{Say it:} \emph{biplôb}
  \item \emph{Say it:} \emph{biplôb}, once more
  \item \emph{Say it:} \emph{biplôber pôre deshe nôtun sôrkār elô.} — After the revolution a new government came to power.
  \item \emph{Recall:} say the Bengali for to surrender, then the Bengali for defeat, then say \emph{biplôb} again
\end{itemize}

\begin{tcolorbox}[breakable,colback=teal!4,colframe=teal!35!black,title={Before you move on}]
What does \bn{বিপ্লব} mean? (\textquotedblleft{}A revolution\textquotedblright{}.) Say it once more.
\end{tcolorbox}
\section[\texorpdfstring{\bn{আন্দোলন} (āndolôn)}{আন্দোলন (āndolôn)}]{\bn{আন্দোলন} (āndolôn) — a movement; a campaign}

\label{lesson:BN-C318-andolon}



\begin{quote}
Before the new one: say the Bengali for defeat, then the Bengali for a revolution.
\end{quote}

\subsection*{You'll want to know: \bn{আন্দোলন}}
\textbf{\bn{আন্দোলন}} — \emph{āndolôn} — \textquotedblleft{}a movement; a campaign\textquotedblright{}.

\begin{quote}
\bn{শহরে} \bn{ছাত্রদের} \bn{একটা} \bn{বড়} \bn{আন্দোলন} \bn{শুরু} \bn{হল।}
\end{quote}

\begin{quote}
\emph{shôhôre chhātrôder ekṭā bôṛo āndolôn shuru hôlô.}
\end{quote}

\begin{quote}
A big student movement began in the city.
\end{quote}

\begin{cousinweb}
\textbf{\bn{আন্দোলন}} (\emph{āndolôn}) is Sanskrit \emph{āndolana}, 'swinging, trembling' (Monier-Williams); the sense 'a political movement', a people set swaying, is the modern one.
\end{cousinweb}

\begin{grammarlens}[title={What you've built}]
One more word for telling what happened.
\end{grammarlens}

\subsection*{Your turn}
\begin{itemize}
\raggedright
  \item \emph{Say it:} \emph{āndolôn}
  \item \emph{Say it:} \emph{āndolôn}, once more
  \item \emph{Say it:} \emph{shôhôre chhātrôder ekṭā bôṛo āndolôn shuru hôlô.} — A big student movement began in the city.
  \item \emph{Recall:} say the Bengali for defeat, then the Bengali for a revolution, then say \emph{āndolôn} again
\end{itemize}

\begin{tcolorbox}[breakable,colback=teal!4,colframe=teal!35!black,title={Before you move on}]
What does \bn{আন্দোলন} mean? (\textquotedblleft{}A movement; a campaign\textquotedblright{}.) Say it once more.
\end{tcolorbox}
\section[\texorpdfstring{\bn{সংগ্রাম} (sônggrām)}{সংগ্রাম (sônggrām)}]{\bn{সংগ্রাম} (sônggrām) — a struggle}

\label{lesson:BN-C318-songgram}



\begin{quote}
Before the new one: say the Bengali for a revolution, then the Bengali for a movement.
\end{quote}

\subsection*{You'll want to know: \bn{সংগ্রাম}}
\textbf{\bn{সংগ্রাম}} — \emph{sônggrām} — \textquotedblleft{}a struggle\textquotedblright{}.

\begin{quote}
\bn{স্বাধীনতার} \bn{জন্য} \bn{সংগ্রাম} \bn{অনেক} \bn{বছর} \bn{চলল।}
\end{quote}

\begin{quote}
\emph{svādhīnôtār jônyô sônggrām ônek bôchhôr chôllô.}
\end{quote}

\begin{quote}
The struggle for independence went on for many years.
\end{quote}

\begin{cousinweb}
\textbf{\bn{সংগ্রাম}} (\emph{sônggrām}) is Sanskrit \emph{saṃgrāma}, first a host of people and then 'battle, conflict' (Monier-Williams).
\end{cousinweb}

\begin{grammarlens}[title={What you've built}]
One more word for telling what happened.
\end{grammarlens}

\subsection*{Your turn}
\begin{itemize}
\raggedright
  \item \emph{Say it:} \emph{sônggrām}
  \item \emph{Say it:} \emph{sônggrām}, once more
  \item \emph{Say it:} \emph{svādhīnôtār jônyô sônggrām ônek bôchhôr chôllô.} — The struggle for independence went on for many years.
  \item \emph{Recall:} say the Bengali for a revolution, then the Bengali for a movement, then say \emph{sônggrām} again
\end{itemize}

\begin{tcolorbox}[breakable,colback=teal!4,colframe=teal!35!black,title={Before you move on}]
What does \bn{সংগ্রাম} mean? (\textquotedblleft{}A struggle\textquotedblright{}.) Say it once more.
\end{tcolorbox}
\section[\texorpdfstring{\bn{প্রতিষ্ঠা} \bn{করা} (prôtishṭhā kôrā)}{প্রতিষ্ঠা করা (prôtishṭhā kôrā)}]{\bn{প্রতিষ্ঠা} \bn{করা} (prôtishṭhā kôrā) — to found; to establish}

\label{lesson:BN-C318-protishtha-kora}



\begin{quote}
Before the new one: say the Bengali for a movement, then the Bengali for a struggle.
\end{quote}

\subsection*{You'll want to know: \bn{প্রতিষ্ঠা} \bn{করা}}
\textbf{\bn{প্রতিষ্ঠা} \bn{করা}} — \emph{prôtishṭhā kôrā} — \textquotedblleft{}to found; to establish\textquotedblright{}.

\begin{quote}
\bn{দাদু} \bn{গ্রামে} \bn{এই} \bn{স্কুলটা} \bn{প্রতিষ্ঠা} \bn{করলেন।}
\end{quote}

\begin{quote}
\emph{dādu grāme ei skulṭā prôtishṭhā kôrlen.}
\end{quote}

\begin{quote}
Grandfather founded this school in the village.
\end{quote}

Past (he/she …): \bn{তিনি} \bn{একটা} \bn{স্কুল} \bn{প্রতিষ্ঠা} \bn{করলেন}, \emph{tini ekṭā skul prôtishṭhā kôrlen}.

\begin{cousinweb}
\textbf{\bn{প্রতিষ্ঠা} \bn{করা}} (\emph{prôtishṭhā kôrā}) puts Sanskrit \emph{pratiṣṭhā}, 'standing still, stability' (Monier-Williams), before \bn{করা}: to make a thing stand firm.
\end{cousinweb}

\begin{grammarlens}[title={What you've built}]
One more word for telling what happened.
\end{grammarlens}

\subsection*{Your turn}
\begin{itemize}
\raggedright
  \item \emph{Say it:} \emph{prôtishṭhā kôrā}
  \item \emph{Say it:} \emph{prôtishṭhā kôrā}, once more
  \item \emph{Say it:} \emph{dādu grāme ei skulṭā prôtishṭhā kôrlen.} — Grandfather founded this school in the village.
  \item \emph{Recall:} say the Bengali for a movement, then the Bengali for a struggle, then say \emph{prôtishṭhā kôrā} again
\end{itemize}

\begin{tcolorbox}[breakable,colback=teal!4,colframe=teal!35!black,title={Before you move on}]
What does \bn{প্রতিষ্ঠা} \bn{করা} mean? (\textquotedblleft{}To found; to establish\textquotedblright{}.) Say it once more.
\end{tcolorbox}
\section[\texorpdfstring{\bn{মুক্তি} (mukti)}{মুক্তি (mukti)}]{\bn{মুক্তি} (mukti) — liberation; release}

\label{lesson:BN-C318-mukti}



\begin{quote}
Before the new one: say the Bengali for a struggle, then the Bengali for to found.
\end{quote}

\subsection*{You'll want to know: \bn{মুক্তি}}
\textbf{\bn{মুক্তি}} — \emph{mukti} — \textquotedblleft{}liberation; release\textquotedblright{}.

\begin{quote}
\bn{যুদ্ধের} \bn{পরে} \bn{কয়েদিরা} \bn{মুক্তি} \bn{পেল।}
\end{quote}

\begin{quote}
\emph{juddher pôre kôyedirā mukti pelô.}
\end{quote}

\begin{quote}
After the war the prisoners were set free.
\end{quote}

\begin{cousinweb}
\textbf{\bn{মুক্তি}} (\emph{mukti}) is Sanskrit \emph{mukti}, 'setting free, release' (Monier-Williams). \bn{মুক্তি} \bn{পাওয়া}, to be set free.
\end{cousinweb}

\begin{grammarlens}[title={What you've built}]
One more word for telling what happened.
\end{grammarlens}

\subsection*{Your turn}
\begin{itemize}
\raggedright
  \item \emph{Say it:} \emph{mukti}
  \item \emph{Say it:} \emph{mukti}, once more
  \item \emph{Say it:} \emph{juddher pôre kôyedirā mukti pelô.} — After the war the prisoners were set free.
  \item \emph{Recall:} say the Bengali for a struggle, then the Bengali for to found, then say \emph{mukti} again
\end{itemize}

\begin{tcolorbox}[breakable,colback=teal!4,colframe=teal!35!black,title={Before you move on}]
What does \bn{মুক্তি} mean? (\textquotedblleft{}Liberation; release\textquotedblright{}.) Say it once more.
\end{tcolorbox}
